\documentclass[a4paper, serif, 10pt]{moderncv}

%% moderncv
% color and style
\moderncvcolor{black}
\moderncvstyle{banking}

%% page layout/margins
\usepackage[top=2.5cm, bottom=2.5cm, left=1.5cm, right=1.5cm]{geometry}

\nopagenumbers{}

%% Babel config for Norwegian language
% ref:
% - https://latex3.github.io/babel/guides/locale-norwegian.html
\usepackage[norsk]{babel}

%% fontspec
\usepackage{fontspec}

\IfFontExistsTF{Linux Libertine}{
  \setmainfont[Ligatures={TeX, Common}, Numbers=OldStyle]{Linux Libertine}
}{}
\IfFontExistsTF{Linux Libertine O}{
  \setmainfont[Ligatures={TeX, Common}, Numbers=OldStyle]{Linux Libertine O}
}{}

%% CTeX
% CJK support, used to show CJK characters in the resume
%
% - heading=false: disable ctex heading
% - punct=kaiming: use kaiming punctuations styles for CJK
% - scheme=plain: use plain scheme, do not override \normalsize font size
% - space=auto: space settings for CJK characters
%
% CJK fonts are resolved per locale language, e.g. Noto Serif/Sans CJK JP
% for Japanese resumes, with fallback to the next region if not installed
%
% ref:
% - http://ctan.mirrorcatalogs.com/language/chinese/ctex/ctex.pdf
\usepackage[UTF8, heading=false, punct=kaiming, scheme=plain, space=auto]{ctex}

\IfFontExistsTF{Noto Serif CJK SC}{
  \setCJKmainfont{Noto Serif CJK SC}
}{}
\IfFontExistsTF{Noto Sans CJK SC}{
  \setCJKsansfont{Noto Sans CJK SC}
}{}

%% line spacing
\usepackage{setspace}
\setstretch{1.125}

%% URL styling - use normal text instead of monospace
\urlstyle{same}

%% Patch \httplink and \httpslink to handle full URLs
% The original moderncv macros always prepend http:// or https:// to URLs.
% This causes issues when the URL already has a protocol (e.g., https://example.com)
% resulting in malformed URLs like https://https://example.com.
% This patch checks if the URL already contains "://" and uses it directly if so.
% Note: We use \str_set:Nx (x-expansion) to expand macros like \@homepage first.
\ExplSyntaxOn
\str_new:N \g_yamlresume_url_str

\RenewDocumentCommand{\httpslink}{O{}m}{
  \str_set:Nx \g_yamlresume_url_str {#2}
  \str_if_in:NnTF \g_yamlresume_url_str {://}
    { \url{#2} }
    { \url{https://#2} }
}

\RenewDocumentCommand{\httplink}{O{}m}{
  \str_set:Nx \g_yamlresume_url_str {#2}
  \str_if_in:NnTF \g_yamlresume_url_str {://}
    { \url{#2} }
    { \url{http://#2} }
}
\ExplSyntaxOff

\name{Silje Bakken}{}
\title{Cybersikkerhetsingeniør | SOC, hendelseshåndtering og skysikkerhet}
\phone[mobile]{+47 123 45 678}
\email{silje.bakken@protonmail.com}
\homepage{https://siljebakken.no}

\address{Storgata 12, Oslo, Oslo, Norge, 0155}{}{}

\extrainfo{{\small \faLinkedin}\ \href{https://www.linkedin.com/in/siljebakken}{@siljebakken} {} {} {} • {} {} {} 
{\small \faGithub}\ \href{https://github.com/siljebakken}{@siljebakken}}

\begin{document}

\maketitle

\section{Grunnleggende informasjon}

\cvline{}{\begin{itemize}
\item Erfaren cybersikkerhetsingeniør med mastergrad i informasjonssikkerhet og 8 års erfaring fra SOC, hendelseshåndtering og sikkerhetsarkitektur.
\item Spesialisert i trusseldeteksjon, logganalyse og automatisering av sikkerhetsprosesser i Microsoft Sentinel, Splunk og QRadar.
\item Har ledet hendelsesrespons for kritisk infrastruktur og redusert responstid gjennom SOAR-playbooks og KQL-baserte deteksjoner.
\item Sterk på risikovurdering, MITRE ATT\&CK, sikkerhetsrevisjoner og samarbeid på tvers av teknologi-, drift- og ledelsesmiljøer.
\item Opptatt av å bygge en proaktiv sikkerhetskultur gjennom opplæring, mentorvirksomhet og tydelig dokumentasjon.
\end{itemize}}
\section{Utdanning}

\cventry{aug. 2014–juni 2016}
        {Master, Informasjonssikkerhet, Poeng: A (5)}
        {Norges teknisk-naturvitenskapelige universitet}
        {\url{https://www.ntnu.no}}
        {}
        {\begin{itemize}
\item Fordypning i kryptografi, nettverkssikkerhet og sikker systemutvikling
\item Gjennomførte masterprosjekt om trusseldeteksjon i kritisk infrastruktur
\item Utviklet praktisk erfaring med risikovurdering, sikkerhetsrevisjon og hendelseshåndtering
\end{itemize}
\textbf{Kurs}: Kryptografi og sikkerhetsprotokoller, Nettsikkerhet, Systemering av sikre systemer, Risikostyring og personvern, Digital etterforskning, Sikkerhetsarkitektur}

\cventry{aug. 2011–juni 2014}
        {Bachelor, Informatikk, Poeng: B (4)}
        {Universitetet i Oslo}
        {\url{https://www.uio.no}}
        {}
        {\begin{itemize}
\item Solid grunnlag i programmering, systemforståelse og datanettverk
\item Arbeidet i prosjektgrupper med sikker utvikling og testing av programvare
\item Skrev bacheloroppgave om autentisering og tilgangskontroll i webløsninger
\end{itemize}
\textbf{Kurs}: Algoritmer og datastrukturer, Operativsystemer, Databaser, Nettverk, Programmering i Java, Diskret matematikk}
\section{Arbeidserfaring}

\cventry{mars 2021–Nå}
        {Senior cybersikkerhetsingeniør}
        {Nordisk Energisikkerhet AS}
        {\url{https://www.nordiskenergisikkerhet.no}}
        {}
        {\begin{itemize}
\item Leder daglig SOC-drift for et energiselskap med kritisk infrastruktur og 24/7-overvåking
\item Utvikler og forvalter deteksjonsregler i Microsoft Sentinel og Splunk med fokus på OT- og IT-konvergens
\item Gjennomfører hendelsesrespons for komplekse hendelser, inkludert analyse, isolering og rapportering
\item Bygget SOAR-playbooks som reduserte manuell prioritering med 35 \% og forkortet gjennomsnittlig responsstid
\item Fungerer som teknisk rådgiver i sikkerhetsarkitektur, risikovurderinger og leverandørrevisjoner
\end{itemize}
\textbf{Nøkkelord}: SOC, Microsoft Sentinel, Splunk, SOAR, OT-sikkerhet, Hendelseshåndtering}

\cventry{aug. 2016–feb. 2021}
        {Cybersikkerhetsingeniør}
        {Bergen Sikkerhetsrådgivning AS}
        {\url{https://www.bergensikkerhet.no}}
        {}
        {\begin{itemize}
\item Overvåket og analyserte sikkerhetshendelser for kunder i finans- og helsesektoren
\item Implementerte og vedlikeholdt SIEM-løsninger (QRadar, ArcSight) og utviklet egne korrelasjonsregler
\item Deltok i hendelsesrespons og digital etterforskning, inkludert innsamling av bevis og rapportering
\item Gjennomførte sårbarhetsanalyser og penetrasjonstester i samarbeid med eksterne konsulenter
\item Bidro til utvikling av sikkerhetspolicyer og opplæring av ansatte i sikker bevissthet
\end{itemize}
\textbf{Nøkkelord}: SIEM, QRadar, ArcSight, Sårbarhetsanalyse, Penetrasjonstesting, Digital etterforskning}
\section{Språk}

\cvline{Norsk}{Morsmål eller tospråklig nivå \hfill \textbf{Nøkkelord}: Morsmål}
\cvline{Engelsk}{Fullt profesjonelt nivå \hfill \textbf{Nøkkelord}: TOEFL 110, IELTS 7.5}
\cvline{Tysk}{Begrenset arbeidskunnskap \hfill \textbf{Nøkkelord}: Grunnleggende}
\section{Ferdigheter}

\cvline{Sikkerhetsovervåking (SOC)}{Ekspert \hfill \textbf{Nøkkelord}: Microsoft Sentinel, Splunk, QRadar, ArcSight, Logganalyse}
\cvline{Trusseldeteksjon og hendelseshåndtering}{Ekspert \hfill \textbf{Nøkkelord}: MITRE ATT\&CK, KQL, SOAR, Playbooks, Digital etterforskning}
\cvline{Skysikkerhet}{Avansert \hfill \textbf{Nøkkelord}: Azure Security, AWS Security, Zero Trust, Identity Management}
\cvline{Risiko og samsvar}{Avansert \hfill \textbf{Nøkkelord}: ISO 27001, GDPR, Risikovurdering, Sikkerhetsrevisjon}
\cvline{Programmering og automatisering}{Viderekommen \hfill \textbf{Nøkkelord}: Python, PowerShell, Bash, KQL, REST API}
\section{Utmerkelser}

\cventry{nov. 2022}
        {Nordisk Energisikkerhet AS}
        {Årets sikkerhetsprosjekt}
        {}
        {}
        {\begin{itemize}
\item Anerkjent for utvikling av en automatisert hendelsesresponsplattform som reduserte responstiden med 40 \%.
\end{itemize}}
\section{Sertifikater}

\cventry{juni 2020}
        {(ISC)²}
        {CISSP}
        {\url{https://www.isc2.org/certifications/cissp}}
        {}
        {}

\cventry{sep. 2021}
        {Microsoft}
        {Azure Security Engineer Associate}
        {\url{https://learn.microsoft.com/en-us/certifications/azure-security-engineer/}}
        {}
        {}

\cventry{mars 2019}
        {ISACA}
        {CISM}
        {\url{https://www.isaca.org/credentialing/cism}}
        {}
        {}
\section{Publikasjoner}

\cventry{okt. 2018}
        {Detecting Advanced Persistent Threats in Critical Infrastructure}
        {Norsk Informasjonssikkerhetskonferanse}
        {\url{https://www.nisk.no/paper/2018}}
        {}
        {\begin{itemize}
\item Presenterer en metode for å oppdage APT-angrep i OT-miljøer ved hjelp av maskinlæring og nettverksanalyse.
\end{itemize}}
\section{Prosjekter}

\cventry{jan. 2020–des. 2020}
        {Et åpent kildekode-prosjekt for å detektere unormal aktivitet i industrielle kontrollsystemer.}
        {Sanntids trusseldeteksjon for OT-nettverk}
        {\url{https://github.com/siljebakken/ot-threat-detection}}
        {}
        {\begin{itemize}
\item Utviklet en deteksjonsmotor basert på Python og Scapy for å analysere Modbus- og DNP3-trafikk.
\item Implementerte varsling til Slack og Microsoft Teams ved mistenkelig aktivitet.
\item Testet løsningen i et simulert OT-miljø med gode resultater.
\end{itemize}
\textbf{Nøkkelord}: OT-sikkerhet, Python, Scapy, Modbus, DNP3}
\section{Interesser}

\cvline{Sport}{Fjellturer, Langrenn, Klatring}
\cvline{Musikk}{Piano, Gitar}
\section{Frivillig arbeid}

\cventry{jan. 2019–Nå}
        {Mentor}
        {Women in Cybersecurity Norge}
        {\url{https://www.wics.no}}
        {}
        {\begin{itemize}
\item Mentorer unge kvinner som ønsker å starte karriere innen cybersikkerhet.
\item Holder workshopserier om SOC-drift og hendelseshåndtering.
\item Bidrar til å øke mangfoldet i sikkerhetsmiljøet.
\end{itemize}}
    
\end{document}